\documentclass[10pt,a4paper]{amsart}
\usepackage{amsmath,amssymb,amsthm,graphicx}
\usepackage{verbatim}
\newenvironment{reference}{\comment}{\endcomment}
%\newenvironment{reference}{}{}
\newenvironment{slogan}{\comment}{\endcomment}
\newenvironment{history}{\comment}{\endcomment}

% For commutative diagrams we use Xy-pic
\usepackage[all]{xy}

% We use 2cell for 2-commutative diagrams.
\xyoption{2cell}
% Theorem environments.
%
\theoremstyle{plain}
\newtheorem{theorem}[subsection]{Theorem}
\newtheorem{proposition}[subsection]{Proposition}
\newtheorem{lemma}[subsection]{Lemma}

\theoremstyle{definition}
\newtheorem{definition}[subsection]{Definition}
\newtheorem{example}[subsection]{Example}
\newtheorem{exercise}[subsection]{Exercise}
\newtheorem{situation}[subsection]{Situation}

\theoremstyle{remark}
\newtheorem{remark}[subsection]{Remark}
\newtheorem{remarks}[subsection]{Remarks}

\numberwithin{equation}{subsection}

% Macros
%
\def\lim{\mathop{\mathrm{lim}}\nolimits}
\def\colim{\mathop{\mathrm{colim}}\nolimits}
\def\Spec{\mathop{\mathrm{Spec}}}
\def\Hom{\mathop{\mathrm{Hom}}\nolimits}
\def\Ext{\mathop{\mathrm{Ext}}\nolimits}
\def\SheafHom{\mathop{\mathcal{H}\!\mathit{om}}\nolimits}
\def\SheafExt{\mathop{\mathcal{E}\!\mathit{xt}}\nolimits}
\def\Sch{\mathit{Sch}}
\def\Mor{\mathop{\mathrm{Mor}}\nolimits}
\def\Ob{\mathop{\mathrm{Ob}}\nolimits}
\def\Sh{\mathop{\mathit{Sh}}\nolimits}
\def\NL{\mathop{N\!L}\nolimits}
\def\CH{\mathop{\mathrm{CH}}\nolimits}
\def\proetale{{pro\text{-}\acute{e}tale}}
\def\etale{{\acute{e}tale}}
\def\QCoh{\mathit{QCoh}}
\def\Ker{\mathop{\mathrm{Ker}}}
\def\Im{\mathop{\mathrm{Im}}}
\def\Coker{\mathop{\mathrm{Coker}}}
\def\Coim{\mathop{\mathrm{Coim}}}

% Boxtimes
%
\DeclareMathSymbol{\boxtimes}{\mathbin}{AMSa}{"02}

%
% Macros for moduli stacks/spaces
%
\def\QCohstack{\mathcal{QC}\!\mathit{oh}}
\def\Cohstack{\mathcal{C}\!\mathit{oh}}
\def\Spacesstack{\mathcal{S}\!\mathit{paces}}
\def\Quotfunctor{\mathrm{Quot}}
\def\Hilbfunctor{\mathrm{Hilb}}
\def\Curvesstack{\mathcal{C}\!\mathit{urves}}
\def\Polarizedstack{\mathcal{P}\!\mathit{olarized}}
\def\Complexesstack{\mathcal{C}\!\mathit{omplexes}}
% \Pic is the operator that assigns to X its picard group, usage \Pic(X)
% \Picardstack_{X/B} denotes the Picard stack of X over B
% \Picardfunctor_{X/B} denotes the Picard functor of X over B
\def\Pic{\mathop{\mathrm{Pic}}\nolimits}
\def\Picardstack{\mathcal{P}\!\mathit{ic}}
\def\Picardfunctor{\mathrm{Pic}}
\def\Deformationcategory{\mathcal{D}\!\mathit{ef}}
\usepackage{cleveref}
\crefname{originallemma}{original lemma}{original lemmas}
\Crefname{originallemma}{Original lemma}{Original lemmas}
\makeatletter
\newlabel{lemma-functor-quasi-coherent-from-affine-diagonal}{{0FZH}{source}{}{sourcepointer-1}{}}
\newlabel{lemma-functor-quasi-coherent-from-affine-diagonal@cref}{{[originallemma][0][]0FZH}{[1][0][]source}{}{}{}}
\makeatother
\crefname{originallemma}{original lemma}{original lemmas}
\Crefname{originallemma}{Original lemma}{Original lemmas}
\makeatletter
\newlabel{lemma-functor-quasi-coherent-from-affine}{{0FZD}{source}{}{sourcepointer-2}{}}
\newlabel{lemma-functor-quasi-coherent-from-affine@cref}{{[originallemma][0][]0FZD}{[1][0][]source}{}{}{}}
\makeatother
\crefname{originalreference}{original reference}{original references}
\Crefname{originalreference}{Original reference}{Original references}
\makeatletter
\newlabel{coherent-lemma-affine-base-change}{{02KG}{source}{}{sourcepointer-3}{}}
\newlabel{coherent-lemma-affine-base-change@cref}{{[originalreference][0][]02KG}{[1][0][]source}{}{}{}}
\makeatother
\crefname{originalremark}{original remark}{original remarks}
\Crefname{originalremark}{Original remark}{Original remarks}
\makeatletter
\newlabel{remark-affine-morphism}{{0FZE}{source}{}{sourcepointer-4}{}}
\newlabel{remark-affine-morphism@cref}{{[originalremark][0][]0FZE}{[1][0][]source}{}{}{}}
\makeatother
\title{Composition of FM functors corresponds to the convolved kernel \$\textbackslash{}mathrm\{pr\}\_\{13,*\}(\textbackslash{}mathrm\{pr\}\_\{12\}\textasciicircum{}*K \textbackslash{}otimes \textbackslash{}mathrm\{pr\}\_\{23\}\textasciicircum{}*L)\$}
\author{The Stacks Project Authors}
\date{stacks-project:functors.tex:local-snapshot-20261012; GNU FDL 1.2 or later}
\begin{document}
\maketitle
\noindent\textit{Editorial note.} All delivered mathematical text is the author's own, copied verbatim from the pinned Stacks Project chapter file functors.tex; no statement or proof step was rewritten, reordered or removed.\par\medskip
\noindent\textit{Editorial note.} Span recomputation by label pairing (LANE\_BRIEF 10.2): in this chapter the author's proof environment follows the statement environment as a standalone block; the spec covers the complete environment pair (statement span, proof span) as recorded in the pool candidate.\par\medskip
\noindent\textit{Editorial note.} Occupation scan: combined claim\_key + statement\_anchor + generic-line scan over all registration files (active and withdrawn) found 0 hits for functors.tex (fu-occupied.json, 2026-10-14); consistent with the duty manager's six-check dispatch.\par\medskip
\ifdefined\hypertarget\hypertarget{sourcepointer-1}{}\fi
\noindent\textit{Source reference.} Original lemma 0FZH: sources/stacks-chow/tags.txt, source lines 19952--19952. Cross-chapter standard prerequisite cited by the author inside the counted span; it is used as a pointer and its own proof is not part of this item.\par\medskip
\ifdefined\hypertarget\hypertarget{sourcepointer-2}{}\fi
\noindent\textit{Source reference.} Original lemma 0FZD: sources/stacks-chow/tags.txt, source lines 19948--19948. Cross-chapter standard prerequisite cited by the author inside the counted span; it is used as a pointer and its own proof is not part of this item.\par\medskip
\ifdefined\hypertarget\hypertarget{sourcepointer-3}{}\fi
\noindent\textit{Source reference.} Original reference 02KG: sources/stacks-chow/tags.txt, source lines 3371--3371. Cross-chapter standard prerequisite cited by the author inside the counted span; it is used as a pointer and its own proof is not part of this item.\par\medskip
\ifdefined\hypertarget\hypertarget{sourcepointer-4}{}\fi
\noindent\textit{Source reference.} Original remark 0FZE: sources/stacks-chow/tags.txt, source lines 19949--19949. Cross-chapter standard prerequisite cited by the author inside the counted span; it is used as a pointer and its own proof is not part of this item.\par\medskip
\begin{lemma}
\label{lemma-compose-FM-QCoh}
Let $R$ be a ring. Let $X$, $Y$, $Z$ be schemes over $R$. Assume
$X$ and $Y$ are quasi-compact and have affine diagonal. Let
$$
F : \QCoh(\mathcal{O}_X) \to \QCoh(\mathcal{O}_Y)
\quad\text{and}\quad
G : \QCoh(\mathcal{O}_Y) \to \QCoh(\mathcal{O}_Z)
$$
be $R$-linear exact functors which commute with arbitrary direct sums.
Let $\mathcal{K}$ in $\QCoh(\mathcal{O}_{X \times_R Y})$
and $\mathcal{L}$ in $\QCoh(\mathcal{O}_{Y \times_R Z})$
be the corresponding ``kernels'', see
Lemma \ref{lemma-functor-quasi-coherent-from-affine-diagonal}.
Then $G \circ F$ corresponds to
$\text{pr}_{13, *}(\text{pr}_{12}^*\mathcal{K}
\otimes_{\mathcal{O}_{X \times_R Y \times_R Z}}
\text{pr}_{23}^*\mathcal{L})$ in $\QCoh(\mathcal{O}_{X \times_R Z})$.
\end{lemma}

\begin{proof}
Since $G \circ F : \QCoh(\mathcal{O}_X) \to \QCoh(\mathcal{O}_Z)$
is $R$-linear, exact, and commutes with arbitrary direct sums,
we find by Lemma \ref{lemma-functor-quasi-coherent-from-affine-diagonal}
that there exists an $\mathcal{M}$ in
$\QCoh(\mathcal{O}_{X \times_R Z})$ corresponding to $G \circ F$.
On the other hand, denote
$\mathcal{E} = \text{pr}_{13, *}(\text{pr}_{12}^*\mathcal{K}
\otimes \text{pr}_{23}^*\mathcal{L})$. Here and in the rest of
the proof we omit the subscript from the tensor products.
Let $U \subset X$ and $W \subset Z$ be affine open subschemes.
To prove the lemma, we will construct an isomorphism
$$
\Gamma(U \times_R W, \mathcal{E})
\cong
\Gamma(U \times_R W, \mathcal{M})
$$
compatible with restriction mappings for varying $U$ and $W$.

\medskip\noindent
First, we observe that
$$
\Gamma(U \times_R W, \mathcal{E}) =
\Gamma(U \times_R Y \times_R W,
\text{pr}_{12}^*\mathcal{K} \otimes \text{pr}_{23}^*\mathcal{L})
$$
by construction. Thus we have to show that the same thing is true
for $\mathcal{M}$.

\medskip\noindent
Write $U = \Spec(A)$ and denote $j : U \to X$ the inclusion morphism.
Recall from the construction of $\mathcal{M}$ in the proof of
Lemma \ref{lemma-functor-quasi-coherent-from-affine} that
$$
\Gamma(U \times_R W, \mathcal{M}) =
\Gamma(W, G(F(j_*\mathcal{O}_U)))
$$
where the $A$-module action on the right hand side is given by
the action of $A$ on $\mathcal{O}_U$. The correspondence between
$F$ and $\mathcal{K}$ tells us that
$F(j_*\mathcal{O}_U) = b_*(a^*j_*\mathcal{O}_U \otimes \mathcal{K})$
where $a : X \times_R Y \to X$ and $b : X \times_R Y \to Y$ are
the projection morphisms. Since $j$ is an affine morphism, we have
$a^*j_*\mathcal{O}_U = (j \times \text{id}_Y)_*\mathcal{O}_{U \times_R Y}$
by Cohomology of Schemes, Lemma
\ref{coherent-lemma-affine-base-change}.
Next, we have
$(j \times \text{id}_Y)_*\mathcal{O}_{U \times_R Y} \otimes \mathcal{K} =
(j \times \text{id}_Y)_*\mathcal{K}|_{U \times_R Y}$
by Remark \ref{remark-affine-morphism} for example.
Putting what we have found together we find
$$
F(j_*\mathcal{O}_U) =
(U \times_R Y \to Y)_*\mathcal{K}|_{U \times_R Y}
$$
with obvious $A$-action. (This formula is implicit in the proof
of Lemma \ref{lemma-functor-quasi-coherent-from-affine}.)
Applying the functor $G$ we obtain
$$
G(F(j_*\mathcal{O}_U)) =
t_*(s^*((U \times_R Y \to Y)_*\mathcal{K}|_{U \times_R Y})
\otimes \mathcal{L})
$$
where $s : Y \times_R Z \to Y$ and $t : Y \times_R Z \to Z$ are
the projection morphisms. Again using affine base change
(Cohomology of Schemes, Lemma \ref{coherent-lemma-affine-base-change})
but this time for the square
$$
\xymatrix{
U \times_R Y \times_R Z \ar[r] \ar[d] & U \times_R Y \ar[d] \\
Y \times_R Z \ar[r] & Y
}
$$
we obtain
$$
s^*((U \times_R Y \to Y)_*\mathcal{K}|_{U \times_R Y}) =
(U \times_R Y \times_R Z \to Y \times_R Z)_*
\text{pr}_{12}^*\mathcal{K}|_{U \times_R Y \times_R Z}
$$
Using Remark \ref{remark-affine-morphism} again we find
\begin{align*}
(U \times_R Y \times_R Z \to Y \times_R Z)_*
\text{pr}_{12}^*\mathcal{K}|_{U \times_R Y \times_R Z}
\otimes \mathcal{L} \\
=
(U \times_R Y \times_R Z \to Y \times_R Z)_*
\left(\text{pr}_{12}^*\mathcal{K} \otimes
\text{pr}_{23}^*\mathcal{L}\right)|_{U \times_R Y \times_R Z}
\end{align*}
Applying the functor $\Gamma(W, t_*(-)) = \Gamma(Y \times_R W, -)$
to this we obtain
\begin{align*}
\Gamma(U \times_R W, \mathcal{M})
& =
\Gamma(W, G(F(j_*\mathcal{O}_U))) \\
& =
\Gamma(Y \times_R W, (U \times_R Y \times_R Z \to Y \times_R Z)_*
(\text{pr}_{12}^*\mathcal{K} \otimes
\text{pr}_{23}^*\mathcal{L})|_{U \times_R Y \times_R Z}) \\
& =
\Gamma(U \times_R Y \times_R W,
\text{pr}_{12}^*\mathcal{K} \otimes \text{pr}_{23}^*\mathcal{L})
\end{align*}
as desired. We omit the verification that these isomorphisms are
compatible with restriction mappings.
\end{proof}

\end{document}
